\documentclass[12pt]{article}
\usepackage[utf8]{inputenc}
\usepackage[T1]{fontenc}
\usepackage[margin=1in]{geometry}
\usepackage{setspace}
\usepackage{xcolor}
\usepackage[colorlinks=true, urlcolor=blue!60!black, linkcolor=blue!60!black]{hyperref}
\setlength{\parindent}{0em}
\setlength{\parskip}{0.6em}
\renewcommand{\baselinestretch}{1.0}
\begin{document}
\begin{center}
{\LARGE \textit{Crosby v. National Foreign Trade Council}, 530 U.S.\ 363 (2000)\par}
\vspace{0.5em}
{\large Case Brief --- Two-Page Edition\par}
\end{center}

\section{Memory Jogger}
Massachusetts's Burma boycott was preempted because it obstructed the federal Burma Act's flexible, calibrated, multilateral sanctions (\href{https://supreme.justia.com/cases/federal/us/530/363/}{Justia}).

\section{Detailed Case Facts}
In 1996 Massachusetts penalized state bidders ``doing business with Burma,'' reaching foreign firms and existing investments, with no waiver or sunset. Months later Congress imposed limited sanctions, authorized the President to ban ``new investment'' by ``United States persons,'' directed a multilateral strategy, and allowed national-security waivers. The EU and Japan protested the state law.

\section{Procedural History}
The National Foreign Trade Council sued; the district court enjoined the law, and the First Circuit affirmed on foreign affairs, Foreign Commerce Clause, and preemption grounds.

\section{Judicial Votes}
Affirmed, 9--0. Souter wrote for seven; Scalia, with Thomas, concurred in the judgment.

\section{Holding}
The state law is preempted as an obstacle to the federal Act's grant of presidential discretion, its limitation to U.S.\ persons and new investment, and its multilateral diplomacy. Other constitutional theories were reserved.

\section{Analysis of Opinions}
\textbf{Majority:} Purchasing power is no defense (\textit{Gould}). The rigid state law reduced the President's ``bargaining chips,'' penalized what Congress chose to permit, and drew foreign protests that impeded diplomacy; congressional silence is ``ambiguous.'' \textbf{Scalia:} the result was ``perfectly obvious on the face of this statute,'' so the legislative history was needless, like St.\ Augustine stealing pears.

\section{Examples: Future Applications}
\textbf{Same-Side:} a broader, unwaivable state ban on firms in a country Congress sanctioned narrowly is preempted. \textbf{Opposite-Side:} a state boycott of a country Congress has not sanctioned falls outside \textit{Crosby}. \textbf{Fence-Sitter:} state pension divestment that tracks a Treasury list and reinvests on delisting.

\section{Critique}
States enjoy a ``market participant'' shield under the dormant Commerce Clause but none against preemption. Treating any partial sanction as a ``calibration of force'' invites expansive preemption, and crediting foreign protests risks a foreign veto over state policy.

\section{Key Quotations}
\begin{quote}
``The fact of a common end hardly neutralizes conflicting means, see \textit{Gade v.\ National Solid Wastes Management Assn.}, 505 U.S.\ 88, 103 (1992), and the fact that some companies may be able to comply with both sets of sanctions does not mean that the state Act is not at odds with achievement of the federal decision about the right degree of pressure to employ. \dots\ Sanctions are drawn not only to bar what they prohibit but to allow what they permit, and the inconsistency of sanctions here undermines the congressional calibration of force.''
\end{quote}
\noindent --- Justice Souter, Opinion of the Court, 530 U.S.\ at 379--380.

\end{document}
